\documentclass[12pt]{article}
% Préambule commun à tous les documents extraits de « La Revue CAPES »
\usepackage[T1]{fontenc}
\usepackage[utf8]{inputenc}
\usepackage{lmodern}
\usepackage{amsmath,amssymb,amsthm,amsfonts,mathrsfs}
\usepackage{setspace}
\usepackage{listings}
\usepackage{pgfplots}
\pgfplotsset{compat=1.18}
\usepackage{longtable}
\usepackage{adjustbox}
\usepackage{lettrine}
\usepackage{tkz-tab}
\usepackage{changepage}
\usepackage{pdflscape}
\usepackage{graphicx}
\usepackage{tikz}
\usepackage{xcolor}
\usepackage[a4paper,top=2cm,bottom=2.5cm,left=2.5cm,right=2cm]{geometry}
\usepackage{fancyhdr}
\usepackage{draftwatermark}
\SetWatermarkText{La RevueCapes}
\SetWatermarkLightness{0.9}
\SetWatermarkScale{2}
\usepackage{hyperref}
\hypersetup{colorlinks=true,linkcolor=blue!50!black,urlcolor=blue!50!black}

\definecolor{revue}{RGB}{20,60,120}

\pagestyle{fancy}
\fancyhf{}
\renewcommand{\headrulewidth}{0.4pt}
\setlength{\headheight}{15pt}

% \entete{Matière}{Titre}{Type de document}
\newcommand{\entete}[3]{%
  \fancyhead[L]{\small\textsc{La Revue CAPES} -- #1}%
  \fancyhead[R]{\small #2}%
  \fancyfoot[C]{\thepage}%
  \thispagestyle{plain}%
  \begin{center}
    {\color{revue}\rule{\textwidth}{1.2pt}}\\[4pt]
    {\small\textsc{Concours directs CAP-CEG / CAPES -- Burkina Faso}}\\[6pt]
    {\LARGE\bfseries\color{revue} #2}\\[6pt]
    {\large\itshape #1 -- #3}\\[2pt]
    {\color{revue}\rule{\textwidth}{1.2pt}}
  \end{center}
  \vspace{0.5cm}%
}

% Encadré pour les remarques ajoutées dans les corrigés
\newcommand{\remarque}[1]{\par\medskip\noindent\fcolorbox{revue}{revue!6}{\parbox{\dimexpr\linewidth-2\fboxsep-2\fboxrule}{\textbf{\color{revue}Remarque.} #1}}\par\medskip}


\begin{document}
\entete{Mathématiques}{Corrigé détaillé du sujet type n°2}{Correction}
\begin{spacing}{1.5}

\textbf{Exercice 1}

(a) Calculons le déterminant de $B$ en fonction de $a,b$ et $c$.


$\det(B)=\begin{vmatrix}
1 & 1 & 1 \\ 
c & a & b \\ 
c^2 & a^2 & b^2
\end{vmatrix}\begin{array}{c}
 c_2\longleftarrow c_2-c_1 \\ 
 c_3\longleftarrow c_3-c_1
 \end{array}  
\begin{vmatrix}
1 & 0 & 0 \\ 
c & a-c & b-c \\ 
c^2 & a^2-c^2 & b^2-c^2
\end{vmatrix} = (a-c)(b-c)\begin{vmatrix}
 1 & 1 \\ 
 a+c & b+c
\end{vmatrix}=(a-c)(b-c)(b-a)$.

(b) Montrer que $\det(A) = 2abc×\det(B)$

$\det(A)=\begin{vmatrix}
a+b & b+c & c+a \\ 
a^2+b^2 & b^2+c^2 & c^2+a^2 \\ 
a^3+b^3 & b^3+c^3 & c^3+a^3
\end{vmatrix}c_1\longleftarrow c_1+c_2+c_3$

$=\begin{vmatrix}
2a+2b+2c & b+c & c+a \\ 
2a^2+2b^2+2c^2 & b^2+c^2 & c^2+a^2 \\ 
2a^3+2b^3+2c^3& b^3+c^3 & c^3+a^3
\end{vmatrix}=2\begin{vmatrix}
a+b+c & b+c & c+a \\ 
a^2+b^2+c^2 & b^2+c^2 & c^2+a^2 \\ 
a^3+b^3+c^3& b^3+c^3 & c^3+a^3
\end{vmatrix}\begin{array}{c}
c_2\longleftarrow c_2-c_1 \\ 
c_3\longleftarrow c_3-c_1
\end{array} $

$=2\begin{vmatrix}
a+b+c & -a & -b \\ 
a^2+b^2+c^2 & -a^2 & -b^2 \\ 
a^3+b^3+c^3& -a^3 &-b^3
\end{vmatrix}c_1\longleftarrow c_1+c_2+c_3=2\begin{vmatrix}
c & -a & -b \\ 
c^2 & -a^2 & -b^2 \\ 
c^3& -a^3 &-b^3
\end{vmatrix}=2abc\begin{vmatrix}
1 & 1 & 1 \\ 
c & a & b \\ 
c^2& a^2 &b^2
\end{vmatrix}=2abc\det(B)$

c) Calculons $\det(A)$

$\det(A)=2abc\det(B)=2abc(a-c)(b-c)(b-a)$

d) Déduisons le rang de $A$ suivant les valeurs de $a, b$ et $c$.


Si $a\neq 0,b\neq 0,c\neq 0,a\neq b,a\neq c,b\neq c$ alors $\det(A)\neq 0$, ainsi $rang(A)=3$

Si $a=0$ ou $b=0$ ou $c=0$ ou $a=b\neq c$ ou $a=c\neq b$ ou $b=c\neq a$, alors $rang(A)=2$

Si $a=b=c\neq 0$ alors $rang(A)=1$

Si $a=b=c=0$ alors $rang(A)=0$

\textbf{Exercice 2}

1. \[A=\begin{pmatrix}
1 & 3 & -3 \\ 
2 & 0 & 2 \\ 
-1 & -1 & 3
\end{pmatrix}\]

2. Déterminons le spectre de $A$

Le polynôme caractéristiques de $A$ est donnée par : 

$P_A(X)=\det(A-XI_3)=\begin{vmatrix}
1-X & 3 & -3 \\ 
2 & -X & 2 \\ 
-1 & -1 & 3-X
\end{vmatrix}\begin{array}{c}
c_1\longleftarrow c_1-c_3\\
c_2\longleftarrow c_2+c_3
\end{array} =\begin{pmatrix}
4-x & 0 & -3 \\ 
0 & 2-x & 2 \\ 
-4+x & 2-x & 3-x
\end{pmatrix}$

\begin{flushright}
$L_2\longleftarrow L_2-L_3$
\end{flushright}

$=(4-x)(2-x)\begin{pmatrix}
1 & 0 & -3 \\ 
1 & 0 & -1+x \\ 
-1 & 1 & 3-x
\end{pmatrix}=-(4-x)(2-x)(2+x)$

Donc $Sp(A)=\{-2;2;4\}$

3. Les sous-espaces propres de $A$

Soit $E_\lambda$ le sous-espace propre associé à la valeur propre $\lambda$.

On a : $E_\lambda=\{u\in \mathbb{R}^3/Au=\lambda u\}$

* Pour $\lambda=-2$, on a:

$\begin{pmatrix}
1 & 3 & -3 \\ 
2 & 0 & 2 \\ 
-1 & -1 & 3
\end{pmatrix}\begin{pmatrix}
x \\ 
y \\ 
z
\end{pmatrix}=-2\begin{pmatrix}
x \\ 
y \\ 
z
\end{pmatrix}\Rightarrow\begin{cases}
x+3y-3z=-2x\\
2x+2z=-2y\\
-x-y+3z=-2z
\end{cases}\Rightarrow u=\begin{pmatrix}
-x \\ 
x \\ 
0
\end{pmatrix}$. Alors $E_{-2}=\langle \begin{pmatrix}
-1 \\ 
1 \\ 
0
\end{pmatrix}\rangle$

* Pour $\lambda=2$, on a :

$\begin{pmatrix}
1 & 3 & -3 \\ 
2 & 0 & 2 \\ 
-1 & -1 & 3
\end{pmatrix}\begin{pmatrix}
x \\ 
y \\ 
z
\end{pmatrix}=2\begin{pmatrix}
x \\ 
y \\ 
z
\end{pmatrix}\Rightarrow\begin{cases}
x+3y-3z=2x\\
2x+2z=2y\\
-x-y+3z=2z
\end{cases}\Rightarrow  u=\begin{pmatrix}
0 \\ 
y \\ 
y
\end{pmatrix}$. Alors $E_{2}=\langle \begin{pmatrix}
0 \\ 
1 \\ 
1
\end{pmatrix}\rangle$

* Pour $\lambda=4$, on a :

$\begin{pmatrix}
1 & 3 & -3 \\ 
2 & 0 & 2 \\ 
-1 & -1 & 3
\end{pmatrix}\begin{pmatrix}
x \\ 
y \\ 
z
\end{pmatrix}=4\begin{pmatrix}
x \\ 
y \\ 
z
\end{pmatrix}\Rightarrow\begin{cases}
x+3y-3z=4x\\
2x+2z=4y\\
-x-y+3z=4z
\end{cases}\Rightarrow u=\begin{pmatrix}
x \\ 
0 \\ 
-x
\end{pmatrix}$. Alors $E_{4}=\langle \begin{pmatrix}
1 \\ 
0 \\ 
-1
\end{pmatrix}\rangle$

4. Diagonalisabilité de $A$

La matrice $A$ est diagonalisable car elle a trois valeurs propres distinctes et trois vecteurs propres linéairement indépendants.Ou autrement car son polynôme caractéristique est scindé en racines simples.

5. Calcul de $A^n$

Étant donné que $A$ est diagonalisable, on peut écrire $A=PDP^{-1}$ avec $P=\begin{pmatrix}
-1 & 0 & 1 \\ 
1 & 1 & 0 \\ 
0 & 1 & -1
\end{pmatrix}$ et $D=\begin{pmatrix}
-2 & 0 & 0 \\ 
0 & 2 & 0 \\ 
0 & 0 & 4
\end{pmatrix}$

Ainsi :$A^n=PD^nP^{-1}$  avec $D^n=\begin{pmatrix}
(-2)^n & 0 & 0 \\ 
0 & 2^n & 0 \\ 
0 & 0 & 4^n
\end{pmatrix}$

Après calcul, $P^{-1}=\frac{1}{2}\begin{pmatrix}
-1 & 1 & -1 \\ 
1 & 1 & 1 \\ 
1 & 1 & -1
\end{pmatrix}$

et $A^n=PD^nP^{-1}=\frac{1}{2}\begin{pmatrix}
-1 & 0 & 1 \\ 
1 & 1 & 0 \\ 
0 & 1 & -1
\end{pmatrix}\begin{pmatrix}
(-2)^n & 0 & 0 \\ 
0 & 2^n & 0 \\ 
0 & 0 & 4^n
\end{pmatrix}\begin{pmatrix}
-1 & 1 & -1 \\ 
1 & 1 & 1 \\ 
1 & 1 & -1
\end{pmatrix}$


$A^n=\frac{1}{2}\begin{pmatrix}
(-2)^n+4^n & -(-2)^n+4^n & (-2)^n-4^n \\ 
-(-2)^n+2^n & (-2)^n+2^n & -(-2)^n+2^n \\ 
2^n-4^n & 2^n-4^n & 2^n+4^n
\end{pmatrix}$

\medskip
6. Calcul de $f\circ f\circ f(u+v+w)=f(f(f(u+v+w)))$

Puisque $f$ est linéaire :

$f(u+v+w)=f(u)+f(v)+f(w)=u+2v-w+3u-w-3u+2v+3w=u+4v+w$.

$f(f(u+v+w))=f(u+4v+w)=f(u)+4f(v)+f(w)=u+2v-w+12u-4w-3u+2v+3w=10u+4v-2w$

$f(f(f(u+v+w)))=f(10u+4v-2w)=10f(u)+4f(v)-2f(w)=28u+16v-20w$

7. Déterminons $a_n,b_n,c_n$ en fonction de $n$.

L'écriture matricielle du système est :


$\begin{pmatrix}
a_{n+1}\\
b_{n+1}\\
c_{n+1}
\end{pmatrix}=\begin{pmatrix}
1 & 3 & -3 \\ 
2 & 0 & 2 \\ 
-1 & -1 & 3
\end{pmatrix}\begin{pmatrix}
a_{n}\\
b_{n}\\
c_{n}
\end{pmatrix}$

Soit $X_n=\begin{pmatrix}
a_{n}\\
b_{n}\\
c_{n}
\end{pmatrix}\Rightarrow X_{n+1}=\begin{pmatrix}
a_{n+1}\\
b_{n+1}\\
c_{n+1}
\end{pmatrix}$

Ainsi $X_{n+1}=AX_n$.

En raisonnant par récurrence, on a: 

$X_n=A^nX_0=\frac{1}{2}\begin{pmatrix}
(-2)^n+4^n & -(-2)^n+4^n & (-2)^n-4^n \\ 
-(-2)^n+2^n & (-2)^n+2^n & -(-2)^n+2^n \\ 
2^n-4^n & 2^n-4^n & 2^n+4^n
\end{pmatrix}\begin{pmatrix}
1\\
1\\
1
\end{pmatrix}$

\[
\Rightarrow \begin{cases}
a_n=\frac{1}{2}((-2)^n+4^n)\\
b_n=\frac{1}{2}(3.2^n-(-2)^n)\\
c_n=\frac{1}{2}(3.2^n-4^n)
\end{cases}
\]


\medskip\textbf{Exercice 3}

1. Tableau des résultats de lancer de 2 dés : 

$
\begin{tabular}{|c|c|c|c|c|c|c|}
\hline 
$X\setminus Y$ & 1 & 2 & 3 & 4 & 5 & 6 \\ 
\hline 
1 & (1;1) & (1;2) & (1;3) & (1;4) & (1;5) & (1;6) \\ 
\hline 
2 & (2;1) & (2;2) & (2;3) & (2;4) & (2;5) & (2;6) \\ 
\hline 
3 & (3;1) & (3;2) & (3;3) & (3;4) & (3;5) & (3;6) \\ 
\hline 
4 & (4;1) & (4;2) & (4;3) & (4;4) & (4;5) & (4;6) \\ 
\hline 
5 & (5;1) & (5;2) & (5;3) & (5;4) & (5;5) & (5;6) \\ 
\hline 
6 & (6;1) & (6;2) & (6;3) & (6;4) & (6;5) & (6;6) \\ 
\hline 
\end{tabular} $

2. La loi de probabilité de la variable aléatoire somme $S=X+Y$,

 Les valeurs possibles de la variables aléatoire  sont :  $\lbrace {2;3;4;5;6;7;8;9;10;11;12 \rbrace}$.
 
 Les variables $X$ et $Y$ sont indépendantes donc
 
$
 P(S=2)=P(X=1)\times P(Y=1)= \frac{1}{6}\times \frac{1}{6}=\frac{1}{36}$ 
 
 $P(S=3)=P(X=1)\times P(Y=2)+P(X=2)\times P(Y=1)= 2\times\frac{1}{6}\times \frac{1}{6}=\frac{2}{36}$

$\cdot\cdot\cdot$
 
On dresse le tableau de la loi de probabilité de la variable aléatoire 

$
\begin{tabular}{|c|c|c|c|c|c|c|c|c|c|c|c|}
\hline 
k & 2 & 3 & 4 & 5 & 6 & 7 & 8 & 9 & 10 & 11 & 12 \\ 
\hline 
P(S=k) & $\frac{1}{36}$ & $\frac{2}{36}$ & $\frac{3}{36}$ & $\frac{4}{36}$ & $\frac{5}{36}$ & $\frac{6}{36}$ & $\frac{5}{36}$ & $\frac{4}{36}$ & $\frac{3}{36}$ & $\frac{2}{36}$ & $\frac{1}{36}$ \\ 
\hline 
\end{tabular}
$

\medskip\textbf{Exercice 4}

1. Soit le système 
 
 $X'=AX$ avec $A=\begin{pmatrix}
 1 & -2 & 0 \\ 
 2 & 0 & -1 \\ 
 4 & -2 & -1
 \end{pmatrix} $
 
 Déterminons les valeurs propres de $A$ puis un système de solutions de $X'=AX$.
 
 Le polynôme caractéristiques de $A$
 
 \begin{align*}
 P_A(x)=\det(A-xI_3) &=\begin{vmatrix}
 1-x & -2 & 0 \\ 
 2 & -x & -1 \\ 
 4 & -2 & -1-x
 \end{vmatrix} L_3\longleftarrow L_3-2L_2\\
 &=\begin{vmatrix}
 1-x & -2 & 0 \\ 
 2 & -x & -1 \\ 
 0 & -2+2x & 1-x
 \end{vmatrix}=(1-x)\begin{vmatrix}
 1-x & -2 & 0 \\ 
 2 & -x & -1 \\ 
 0 & -2 & 1
 \end{vmatrix}L_2\longleftarrow L_2+L_3\\
 &=(1-x)\begin{vmatrix}
 1-x & -2 & 0 \\ 
 2 & -2-x & 0 \\ 
 0 & -2 & 1
 \end{vmatrix}\\
 &=(1-x)((1-x)(-2-x)+4)=(1-x)(x^2+x+2)
 \end{align*}
 
$\Delta=1-8=-7=(i\sqrt{7})^2\Rightarrow x_1=\frac{-1-i\sqrt{7}}{2}$ et $x_2=\frac{-1+i\sqrt{7}}{2}$

Donc, $P_A(x)=(1-x)(x+\frac{1+i\sqrt{7}}{2})(x+\frac{1-i\sqrt{7}}{2})$

Ainsi les valeurs propres de $A$ : $Sp(A)=\{1;\frac{-1-i\sqrt{7}}{2};\frac{-1+i\sqrt{7}}{2}\}$

\medskip
Calculons les sous-espaces propres  et vecteurs propres :

-$E_1=\{u=\begin{pmatrix}
x\\
y\\
z
\end{pmatrix}\in \mathbb{R}^3/Au=u\}$

$Au=u\Rightarrow \begin{pmatrix}
 1 & -2 & 0 \\ 
 2 & 0 & -1 \\ 
 4 & -2 & -1
 \end{pmatrix}\begin{pmatrix}
x\\
y\\
z
\end{pmatrix}=\begin{pmatrix}
x\\
y\\
z
\end{pmatrix}\Rightarrow \begin{cases}
x-2y=x\\
2x-z=y\\
4x-2y-z=z
\end{cases}\Rightarrow \begin{cases}
y=0\\
z=2x
\end{cases}$

$\Rightarrow u=\begin{pmatrix}
x\\
0\\
2x
\end{pmatrix}=x\begin{pmatrix}
1\\
0\\
2
\end{pmatrix}\Rightarrow E_1=\langle\begin{pmatrix}
1\\
0\\
2
\end{pmatrix}\rangle$

\medskip
-$E_{\frac{-1-i\sqrt{7}}{2}}=\{v=\begin{pmatrix}
x\\
y\\
z
\end{pmatrix}\in \mathbb{C}^3/Av=\frac{-1-i\sqrt{7}}{2}v\}$

\begin{align*}
Av=\frac{-1-i\sqrt{7}}{2}v &\Rightarrow \begin{cases}
x-2y=\frac{-1-i\sqrt{7}}{2}x\\
2x-z=\frac{-1-i\sqrt{7}}{2}y\\
4x-2y-z=\frac{-1-i\sqrt{7}}{2}z
\end{cases}\Rightarrow\begin{cases}
2x-4y=-x-i\sqrt{7}x\\
4x-2z=-y-i\sqrt{7}y\\
8x-4y-2z=-z-i\sqrt{7}z
\end{cases}\\
&\Rightarrow \begin{cases}
(3+i\sqrt{7})x-4y=0\\
4x+(1+i\sqrt{7})y-2z=0\\
8x-4y-(1-i\sqrt{7})z=0
\end{cases}
\end{align*}

$L_1\Longrightarrow y=\frac{3+i\sqrt{7}}{4}x$ (*)

(*) dans $L_3\Longrightarrow z=\frac{3+i\sqrt{7}}{2}x$

$\Rightarrow v=\begin{pmatrix}
x\\
\frac{3+i\sqrt{7}}{4}x\\
\frac{3+i\sqrt{7}}{2}x
\end{pmatrix}= \frac{3+i\sqrt{7}}{4}x\begin{pmatrix}
 \frac{3-i\sqrt{7}}{4}\\
1\\
2
\end{pmatrix}\Rightarrow E_{\frac{-1-i\sqrt{7}}{2}}=\langle\begin{pmatrix}
\frac{3-i\sqrt{7}}{4}\\
1\\
2
\end{pmatrix}\rangle$

\medskip
-$E_{\frac{-1+i\sqrt{7}}{2}}=\{w=\begin{pmatrix}
x\\
y\\
z
\end{pmatrix}\in \mathbb{C}^3/Aw=\frac{-1+i\sqrt{7}}{2}w\}$

\begin{align*}
Aw=\frac{-1+i\sqrt{7}}{2}w &\Rightarrow \begin{cases}
x-2y=\frac{-1+i\sqrt{7}}{2}x\\
2x-z=\frac{-1+i\sqrt{7}}{2}y\\
4x-2y-z=\frac{-1+i\sqrt{7}}{2}z
\end{cases}\Rightarrow\begin{cases}
2x-4y=-x+i\sqrt{7}x\\
4x-2z=-y+i\sqrt{7}y\\
8x-4y-2z=-z+i\sqrt{7}z
\end{cases}\\
&\Rightarrow \begin{cases}
(3-i\sqrt{7})x-4y=0\\
4x+(1-i\sqrt{7})y-2z=0\\
8x-4y-(1+i\sqrt{7})z=0
\end{cases}
\end{align*}

$L_1\Longrightarrow y=\frac{3-i\sqrt{7}}{4}x$ (*)

(*) dans $L_3\Longrightarrow z=\frac{3-i\sqrt{7}}{2}x$

$\Rightarrow v=\begin{pmatrix}
x\\
\frac{3-i\sqrt{7}}{4}x\\
\frac{3-i\sqrt{7}}{2}x
\end{pmatrix}= \frac{3-i\sqrt{7}}{4}x\begin{pmatrix}
 \frac{3+i\sqrt{7}}{4}\\
1\\
2
\end{pmatrix}\Rightarrow E_{\frac{-1+i\sqrt{7}}{2}}=\langle\begin{pmatrix}
\frac{3+i\sqrt{7}}{4}\\
1\\
2
\end{pmatrix}\rangle$

Donc, la matrice diagonale $D=\begin{pmatrix}
1 & 0 & 0 \\ 
0 & \frac{-1-i\sqrt{7}}{2} & 0 \\ 
0 & 0 & \frac{-1+i\sqrt{7}}{2}
\end{pmatrix} $ et la matrice de passage est $P=\begin{pmatrix}
1 & \frac{3-i\sqrt{7}}{4} & \frac{3+i\sqrt{7}}{4} \\ 
0 & 1                   & 1 \\
2 & 2                     & 2 
\end{pmatrix}$

Calculons $P^{-1}$

$
\det(P) =\begin{vmatrix}
1 & \frac{3-i\sqrt{7}}{4} & \frac{3+i\sqrt{7}}{4} \\ 
0 & 1                   & 1 \\
2 & 2                     & 2 
\end{vmatrix}c_3\longleftarrow c_3-c_2=\begin{vmatrix}
1 & \frac{3-i\sqrt{7}}{4} & \frac{i\sqrt{7}}{2} \\ 
0 & 1                   & 0 \\
2 & 2                     & 0 
\end{vmatrix}=-2\times\frac{i\sqrt{7}}{2}=-i\sqrt{7}
$

$P^{-1}=\frac{1}{\det(P)}t_{com(P)}$

$com(P)=\begin{pmatrix}
0 & 2 & -2 \\ 
i\sqrt{7} & \frac{1-i\sqrt{7}}{2} & -\frac{1+i\sqrt{7}}{2} \\ 
\frac{-i\sqrt{7}}{2} & -1 & 1
\end{pmatrix} \Rightarrow t_{com(P)}=\begin{pmatrix}
0 & i\sqrt{7} & \frac{-i\sqrt{7}}{2} \\ 
2 & \frac{1-i\sqrt{7}}{2} & -1 \\ 
-2 & -\frac{1+i\sqrt{7}}{2} & 1
\end{pmatrix}$

Donc, $P^{-1}=-\frac{1}{i\sqrt{7}}\begin{pmatrix}
0 & i\sqrt{7} & \frac{-i\sqrt{7}}{2} \\ 
2 & \frac{1-i\sqrt{7}}{2} & -1 \\ 
-2 & -\frac{1+i\sqrt{7}}{2} & 1
\end{pmatrix}=\frac{i\sqrt{7}}{7}\begin{pmatrix}
0 & i\sqrt{7} & \frac{-i\sqrt{7}}{2} \\ 
2 & \frac{1-i\sqrt{7}}{2} & -1 \\ 
-2 & -\frac{1+i\sqrt{7}}{2} & 1
\end{pmatrix}$

Ainsi $A=PDP^{-1}$

\medskip
L'équation différentielle $X'=AX$, admet une solution de la forme,

$X(t)=e^{At}X(0)$ avec ,$e^{At}=Pe^{Dt}P^{-1}$, d'où $X(t)=Pe^{Dt}P^{-1}X(0)$ 

$e^{Dt}=\begin{pmatrix}
e^t & 0 & 0 \\ 
0 & e^{\frac{-1-i\sqrt{7}}{2}t} & 0 \\ 
0 & 0 & e^{\frac{-1+i\sqrt{7}}{2}t}
\end{pmatrix}$


\newgeometry{left=1cm, right=1cm, top=1cm, bottom=1cm}% Suppression des marges

\resizebox{\textwidth}{!}{
$X(t)=\begin{pmatrix}
x(t)\\
y(t)\\
z(t)
\end{pmatrix}=\frac{i\sqrt{7}}{7}\begin{pmatrix}
1 & \frac{3-i\sqrt{7}}{4} & \frac{3+i\sqrt{7}}{4} \\ 
0 & 1                   & 1 \\
2 & 2                     & 2 
\end{pmatrix}\begin{pmatrix}
e^t & 0 & 0 \\ 
0 & e^{\frac{-1-i\sqrt{7}}{2}t} & 0 \\ 
0 & 0 & e^{\frac{-1+i\sqrt{7}}{2}t}
\end{pmatrix}\begin{pmatrix}
0 & i\sqrt{7} & \frac{-i\sqrt{7}}{2} \\ 
2 & \frac{1-i\sqrt{7}}{2} & -1 \\ 
-2 & -\frac{1+i\sqrt{7}}{2} & 1
\end{pmatrix}\begin{pmatrix}
x(0)\\
y(0)\\
z(0)
\end{pmatrix}$}


\resizebox{\textwidth}{!}{$\begin{pmatrix}
x(t)\\
y(t)\\
z(t)
\end{pmatrix}=\frac{i\sqrt{7}}{7}\begin{pmatrix}
\frac{3-i\sqrt{7}}{2}e^{\frac{-1-i\sqrt{7}}{2}t} -\frac{3+i\sqrt{7}}{2}e^{\frac{-1+i\sqrt{7}}{2}t}& i\sqrt{7}e^t- \frac{1+i\sqrt{7}}{2}e^{\frac{-1-i\sqrt{7}}{2}t}-\frac{1-i\sqrt{7}}{2}e^{\frac{-1+i\sqrt{7}}{2}t}& \frac{-i\sqrt{7}}{2}e^t-\frac{3-i\sqrt{7}}{4}e^{\frac{-1-i\sqrt{7}}{2}t}+\frac{3+i\sqrt{7}}{4}e^{\frac{-1+i\sqrt{7}}{2}t} \\ 
2e^{\frac{-1-i\sqrt{7}}{2}t} -2e^{\frac{-1+i\sqrt{7}}{2}t}& \frac{1-i\sqrt{7}}{2}e^{\frac{-1-i\sqrt{7}}{2}t}-\frac{1+i\sqrt{7}}{2}e^{\frac{-1+i\sqrt{7}}{2}t} & -e^{\frac{-1-i\sqrt{7}}{2}t}+e^{\frac{-1+i\sqrt{7}}{2}t} \\ 
4e^{\frac{-1-i\sqrt{7}}{2}t}-4e^{\frac{-1+i\sqrt{7}}{2}t} & 2i\sqrt{7}e^t+(1-i\sqrt{7})e^{\frac{-1-i\sqrt{7}}{2}t} -(1+i\sqrt{7})e^{\frac{-1+i\sqrt{7}}{2}t} & -i\sqrt{7}e^t-2e^{\frac{-1-i\sqrt{7}}{2}t}+2e^{\frac{-1+i\sqrt{7}}{2}t}
\end{pmatrix}\begin{pmatrix}
x(0)\\
y(0)\\
z(0)
\end{pmatrix} $
}

\medskip




\resizebox{\textwidth}{!}{$\begin{cases}
x(t)=\frac{i\sqrt{7}}{7}[(\frac{3-i\sqrt{7}}{2}e^{\frac{-1-i\sqrt{7}}{2}t} -\frac{3+i\sqrt{7}}{2}e^{\frac{-1+i\sqrt{7}}{2}t})x(0)+(i\sqrt{7}e^t- \frac{1+i\sqrt{7}}{2}e^{\frac{-1-i\sqrt{7}}{2}t}-\frac{1-i\sqrt{7}}{2}e^{\frac{-1+i\sqrt{7}}{2}t})y(0)+(\frac{-i\sqrt{7}}{2}e^t-\frac{3-i\sqrt{7}}{4}e^{\frac{-1-i\sqrt{7}}{2}t}+\frac{3+i\sqrt{7}}{4}e^{\frac{-1+i\sqrt{7}}{2}t})]\\
y(t)=\frac{i\sqrt{7}}{7}[(2e^{\frac{-1-i\sqrt{7}}{2}t} -2e^{\frac{-1+i\sqrt{7}}{2}t})x(0)+(\frac{1-i\sqrt{7}}{2}e^{\frac{-1-i\sqrt{7}}{2}t}-\frac{1+i\sqrt{7}}{2}e^{\frac{-1+i\sqrt{7}}{2}t})y(0)+(-e^{\frac{-1-i\sqrt{7}}{2}t}+e^{\frac{-1+i\sqrt{7}}{2}t})z(0)]\\
z(t)=\frac{i\sqrt{7}}{7}[(4e^{\frac{-1-i\sqrt{7}}{2}t}-4e^{\frac{-1+i\sqrt{7}}{2}t})x(0)+(2i\sqrt{7}e^t+(1-i\sqrt{7})e^{\frac{-1-i\sqrt{7}}{2}t} -(1+i\sqrt{7})e^{\frac{-1+i\sqrt{7}}{2}t})y(0)+(-i\sqrt{7}e^t-2e^{\frac{-1-i\sqrt{7}}{2}t}+2e^{\frac{-1+i\sqrt{7}}{2}t})z(0)]
\end{cases}$
}

\restoregeometry %Restaurercles mages normales



Avec $x(0),y(0)$ et $z(0)\in \mathbb{R}$ déterminés par les conditions initiales.

\textbf{\underline{N.B :}} La méthode utilisez ici est un peu plus longue. On pourrait utilisez d'autre méthode plus courte, qui consisterait à écrire pour chaque valeur propre une solution fondamentale de la forme $x_i(t)=e^{\lambda t}u$ ou $\lambda$ est la valeur propre et $u$ le vecteur propre associé à la valeur propre $\lambda$.

Ensuite la solution est donnée par $X(t)=x(0)x_1(t)+y(0)x_2(t)+z(0)x_3(t)$.

\medskip
2. Pour :

$\bullet A=\begin{pmatrix}
1 & -1 \\
4 & -3
\end{pmatrix}$

Le polynôme caractéristique de $A$

$P_A(x)=\det(A-xI_2)=\begin{vmatrix}
1-x & -1 \\
4 & -3-x
\end{vmatrix}=(1-x)(-3-x)+4=x^2+2x+1=(x+1)^2$

$A$ a une seule valeur propre -1.

Calculons le(s) vecteur(s) propre(s) :

Soit $u=\begin{pmatrix}
x\\
y
\end{pmatrix}\in \mathbb{R}^2$ le vecteur propre associé à la valeur propre -1.

$Au=-u\Rightarrow \begin{pmatrix}
1 & -1 \\
4 & -3
\end{pmatrix}\begin{pmatrix}
x\\
y
\end{pmatrix}=-\begin{pmatrix}
x\\
y
\end{pmatrix}\Rightarrow \begin{cases}
x-y=-x\\
4x-3y=-y
\end{cases}\Rightarrow y=2x$

$u=\begin{pmatrix}
x\\
y
\end{pmatrix}=\begin{pmatrix}
x\\
2x
\end{pmatrix}=x\begin{pmatrix}
1\\
2
\end{pmatrix}\Rightarrow u=\begin{pmatrix}
1\\
2
\end{pmatrix}$ et $E_{-1}=\langle\begin{pmatrix}
1\\
2
\end{pmatrix}\rangle$

$A$ n'est pas diagonalisable car la valeur propre -1 de multiplicité 2 à une seule valeur propre.

Solution de l'équation $X'=AX$

Étant donné que cette matrice d'ordre 2 a une seule valeur propre de multiplicité 2, à laquelle est associé un seule vecteur propre $u$ (ce qui est insuffisant pour une base complète), elle est dite \textbf{défectueuse}. Cela signifie que nous devons utiliser un vecteur propre généralisé pour obtenir une solution complète. 

Déterminons cet vecteur pour compléter la base.

Soit $v$ cet vecteur, il est obtenir par la formule :

$(A+I)v=u\Longrightarrow \begin{pmatrix}
2 & -1 \\
4 & -2
\end{pmatrix}\begin{pmatrix}
x\\
y
\end{pmatrix}=\begin{pmatrix}
1\\
2
\end{pmatrix}\Rightarrow \begin{cases}
2x-y=1\\
4x-2y=2
\end{cases}
\Rightarrow 2x-y=1\Rightarrow y=2x-1$

Soit $v=\begin{pmatrix}
0\\
-1
\end{pmatrix}$ (qui vérifie $y=2x-1$)

Vérifions que $u$ et $v$ sont libre.

On $\det(\begin{pmatrix}
1 & 0\\
2 & -1
\end{pmatrix})=\begin{vmatrix}
1 & 0\\
2 & -1
\end{vmatrix}=-1\neq 0$ alors $u$ et $v$ sont libres. ils forment donc une base de $\mathbb{R}^2$.

\medskip
\textbf{Solution du système}

La solution associée au vecteur propre est $X_1(t)=e^{-t}\begin{pmatrix}
1\\
2
\end{pmatrix}=\begin{pmatrix}
e^{-t}\\
2e^{-t}
\end{pmatrix}$


La solution associée au vecteur propre généralisé est $X_2(t)=te^{-t}\begin{pmatrix}
1\\
2
\end{pmatrix} + e^{-t}\begin{pmatrix}
0\\
-1
\end{pmatrix}$

L'équation différentielle $X'=AX$, admet une solution de la forme,

\resizebox{\textwidth}{!}{
$X(t)=\begin{pmatrix}
x(t)\\
y(t)
\end{pmatrix}=x(0)X_1(t)+y(0)X_2(t)= x(0)\begin{pmatrix}
e^{-t}\\
2e^{-t}
\end{pmatrix}+y(0)\left( te^{-t}\begin{pmatrix}
1\\
2
\end{pmatrix} + e^{-t}\begin{pmatrix}
0\\
-1
\end{pmatrix}\right)=\begin{pmatrix}
x(0)e^{-t}+y(0)te^{-t}\\
2x(0)e^{-t}+2y(0)te^{-t}-y(0)e^{-t}
\end{pmatrix}$}


$\Rightarrow\begin{cases}
x(t)=x(0)e^{-t}+y(0)te^{-t}\\
y(t)=(2x(0)-y(0))e^{-t}+2y(0)te^{-t}
\end{cases}$


$\bullet$ On poursuivra les mêmes procédures pour la matrice $A=\begin{pmatrix}
 2 & 1 & 0 & 0 \\ 
 0 & 0 & -1 & 0 \\ 
 1 & 2 & 1 & 0 \\ 
 -1 & 2 & 2 & 1
 \end{pmatrix} $.

\end{spacing}
\end{document}
